\chapter{Podsumowanie}

\section{Ocena stopnia realizacji założeń projektowych}

Projekt zrealizował główne założenia dla tematu \textit{System zarządzania hotelem}. Aplikacja pozwala na podstawowe zarządzanie gośćmi, pokojami, rezerwacjami, płatnościami, pracownikami i usługami dodatkowymi. Został wykonany interfejs graficzny, lokalna baza danych SQLite, migracje Entity Framework Core oraz prosta walidacja danych.

Program nie jest pełnym systemem hotelowym klasy komercyjnej, ale nie taki był główny zakres projektu. Dla podstawowego systemu zarządzania wykonano najważniejsze funkcje administracyjne, które pozwalają kontrolować dane recepcji hotelowej.

\section{Napotkane problemy}

Podczas realizacji projektu mogły pojawić się typowe problemy związane z połączeniem aplikacji WPF z bazą danych. Jednym z problemów była poprawna konfiguracja Entity Framework Core i SQLite, aby baza tworzyła się automatycznie przy uruchomieniu programu.

Drugim problemem było sprawdzanie dostępności pokoi. Trzeba było uwzględnić sytuację, w której dwie rezerwacje nakładają się na siebie czasowo. W tym celu w serwisie sprawdzany jest konflikt dat dla wybranego pokoju.

Kolejnym problemem było ustawienie interfejsu tak, aby panel główny nie wyglądał zbyt pusto. Dlatego dodano karty statystyk, opis modułów oraz poprawiono wygląd zakładki usług.

\section{Możliwe kierunki dalszego rozwoju aplikacji}

Projekt można rozbudować do pełniejszego systemu dla hotelu. Można dodać logowanie i role użytkowników, na przykład administratora, recepcjonisty i menedżera. Można też dodać faktury, raporty miesięczne, eksport do pliku PDF, historię zmian oraz pełne zarządzanie pokojami.

W przyszłości można również dodać moduł rezerwacji online albo integrację z zewnętrznymi portalami. Obecna wersja jest jednak podstawowym systemem zarządzania i realizuje zakres przyjęty w projekcie.
